\LessonHeader{2}{China's Fujian Throws Jets into the Sky}{中国福建舰把战机抛上天空}{科学新闻 · Science News}{话题\ 航空母舰·电磁弹射\enspace|\enspace 新词\ \textbf{103}\enspace|\enspace 新义\ \textbf{14}\enspace|\enspace 短语\ \textbf{26}\enspace|\enspace 全文\ \textbf{498}\ 词}{Lesson 2\enspace·\enspace China's Fujian Throws Jets into the Sky}
\begin{lessonbody}
\PARA{1}\SMK{g}On 22 \NW{}{September}{,}{/sepˈtembə(r)/}{n.}{九月} \NW{}{China}{}{/ˈtʃaɪnə/}{n.}{中国} \NWB{}{shared}{}{share}{/ʃeə(r)/}{vt.}{分享；公布} some big \NW{}{news}{.}{/njuːz/}{n.[U]}{新闻；消息}\EMK{g}{1} 
\GRAM{1}{\GITEM{一般过去时}{一般过去时表示过去某个时间发生的事，事情已经结束。}{shared 是 share 加 -d，说 9 月 22 日那天发生的事。}}{}
\SMK{g}Three \PH{\NWB{}{kinds}{}{kind}{/kaɪnd/}{n.[C]}{种类}\ of}{kind(s) of}{……种（类）} \DV{}{carrier}{}{carry + -er → 名词}{n.[C]}{航空母舰} \NWB{}{planes}{,}{plane}{/pleɪn/}{n.[C]}{飞机} the \PN{}{J-15T}{}{/ˌdʒeɪ ˌfɪfˈtiːn ˈtiː/}{n.}{歼-15T} and \PN{}{J-35}{}{/ˌdʒeɪ ˌθɜːti ˈfaɪv/}{n.}{歼-35} \DV{}{fighters}{}{fight + -er → 名词}{n.[C]}{战斗机} and the \PN{}{KJ-600}{}{/ˌkeɪ ˌdʒeɪ sɪks ˈhʌndrəd/}{n.}{空警-600} \NW{}{early}{-}{/ˈɜːli/}{adj.}{早的；提前的}\NW{}{warning}{}{/ˈwɔːnɪŋ/}{n.[U]}{警告；预警} \AG{}{plane}{,} had \NWB{}{finished}{}{finish}{/ˈfɪnɪʃ/}{vt.}{完成} their first \NW{}{training}{}{/ˈtreɪnɪŋ/}{n.[U]}{训练} on the \AG{}{carrier}{} \PN{}{Fujian}{.}{}{n.}{福建舰}\EMK{g}{2} 
\GRAM{2}{\GITEM{过去完成时}{过去完成时表示在过去某个时间或某个动作之前已经完成的事。}{had finished：宣布之前训练已经完成。finished 与过去式同形。}\GITEM{复合形容词}{用连字符连成的形容词放在名词前，其中名词用单数。}{early-warning 两词连成形容词，放在 plane 前，指“预警”飞机。}}{}
\SMK{g}An \DV{}{electromagnetic}{}{electro- + magnetic（磁的）}{adj.}{电磁的} \GL{}{catapult}{}{/ˈkætəpʌlt/}{n.[C]}{弹射器} \NWB{}{threw}{}{throw}{/θrəʊ/}{vt.}{抛；扔} them into the \NS{}{air}{,}{新义}{}{/eə(r)/}{n.[U]}{空中} \RV{}{even}{} when they were \PH{\NW{}{full}{}{/fʊl/}{adj.}{满的；装满的}\ of}{be full of}{装满；充满} \NW{}{oil}{}{/ɔɪl/}{n.[U]}{油；燃油} and \GL{}{weapons}{.}{/ˈwepən/}{n.[C]}{武器}\EMK{g}{3}
\GRAM{3}{\GITEM{时间状语从句}{时间状语从句在句中作时间状语，说明主句动作发生的时间。}{when they were full of… 说弹射是在什么时候；even when 意为“即使在……时”。}}{\SEG{An electromagnetic catapult}{主语（名词短语）}{segA}\SEG{threw}{谓语（动词过去式）}{segB}\SEG{them}{宾语（代词）}{segC}\SEG{into the air}{状语（介词短语，说去向）}{segD}\SEG{even when they were full of oil and weapons}{时间状语从句（说何时）}{segA}}

\PARA{2}\PH{\NWB{}{Hats}{}{hat}{/hæt/}{n.[C]}{帽子}\ \NW{}{off}{}{/ɒf/}{adv.}{脱下}\ to}{hats off (to sb)}{向……致敬} that! \SMK{g}\NW{}{Ah}{,}{/ɑː/}{int.}{啊} but what \NW{}{exactly}{}{/ɪɡˈzæktli/}{adv.}{确切地；到底} is an \AG{}{electromagnetic}{} catapult, and why is it \NW{}{such}{}{/sʌtʃ/}{det.}{这样的} \PH{a\ big\ \NW{}{deal}{?}{/diːl/}{n.[C]}{大事；事情}}{a big deal}{了不起的事}\EMK{g}{4}
\GRAM{4}{\GITEM{疑问句}{一般疑问句问“是不是”，特殊疑问句问具体信息，反意疑问句请对方确认。}{what、why 开头问具体情况；is 放到主语 it 前面。}\GITEM{感叹词和语气}{oh、wow、please 等表示语气，不作句子成分。}{Ah 单独用，表示“想到一个问题”的语气。}}{}

\PARA{3}A \AG{}{carrier}{} is like an \NW{}{airport}{}{/ˈeəpɔːt/}{n.[C]}{机场} \PH{at\ \NW{}{sea}{.}{/siː/}{n.[U]}{海；海上}}{at sea}{在海上} \SMK{g}A \AG{}{plane}{} has to \PH{\NW{}{speed}{}{/spiːd/}{vi.}{加速}\ up}{speed up}{加速} before it can \NW{}{rise}{}{/raɪz/}{vi.}{升起；上升} into the \AG{}{air}{.}\EMK{g}{5} 
\GRAM{5}{\GITEM{have to}{have to do 由于客观情况“不得不”；don't have to 不必。}{has to speed up：飞机为了起飞“不得不”先加速，has to 后接原形。}\GITEM{情态动词}{情态动词后接动词原形，表示能力、可能、许可、义务或建议等。}{can rise：can 后接原形 rise，说飞机“能”升空。}}{}
But today’s \AG{}{carrier}{} \AG{}{planes}{} \NW{}{weigh}{}{/weɪ/}{vt.}{有……重} \PH{tens\ of}{tens of}{几十} \NWB{}{tons}{,}{ton}{/tʌn/}{n.[C]}{吨} and a \GL{}{deck}{}{/dek/}{n.[C]}{甲板} of a \RV{}{few}{} \NW{}{hundred}{}{/ˈhʌndrəd/}{num.}{百} \NWB{}{metres}{}{metre}{/ˈmiːtə(r)/}{n.[C]}{米} is no \RVB{}{longer}{}{long} \NW{}{enough}{.}{/ɪˈnʌf/}{adj.}{足够的}

\PARA{4}The \NW{}{usual}{}{/ˈjuːʒuəl/}{adj.}{通常的} \NW{}{answer}{}{/ˈɑːnsə(r)/}{n.[C]}{答案；办法} is the \GL{}{steam}{}{/stiːm/}{n.[U]}{蒸汽} catapult. Steam \NW{}{pressure}{}{/ˈpreʃə(r)/}{n.[U]}{压力} \AGB{}{throws}{}{throw} the \AG{}{plane}{} \NW{}{forward}{,}{/ˈfɔːwəd/}{adv.}{向前} but the \RV{}{system}{} is \NW{}{slow}{.}{/sləʊ/}{adj.}{慢的} \PH{In\ an\ \NW{}{emergency}{,}{/iˈmɜːdʒənsi/}{n.[C]}{紧急情况}}{in an emergency}{在紧急情况下} the \NW{}{ship}{}{/ʃɪp/}{n.[C]}{舰；船} first has to \PH{\NW{}{bring}{}{/brɪŋ/}{vt.}{使达到}\ water\ to\ the\ \NW{}{boil}{}{/bɔɪl/}{n.[sing.]}{沸腾}}{bring sth to the boil}{把……烧开} and \NW{}{wait}{}{/weɪt/}{vi.}{等待} for steam. So \NW{}{could}{}{/kʊd/}{modal v.}{能；可能} \NW{}{electricity}{}{/ɪˌlekˈtrɪsəti/}{n.[U]}{电} do the \NW{}{job}{}{/dʒɒb/}{n.[C]}{活儿；任务} \NW{}{quickly}{?}{/ˈkwɪkli/}{adv.}{快地；迅速地}

\PARA{5}\SMK{g}\AG{}{Electricity}{} \NS{}{sounds}{}{新词性}{sound}{/saʊnd/}{vi.}{听起来} \NW{}{easy}{}{/ˈiːzi/}{adj.}{容易的} to \NS{}{use}{:}{新词性}{}{/juːz/}{vt.}{使用} you \NW{}{press}{}{/pres/}{vt.}{按} a \NW{}{button}{,}{/ˈbʌtn/}{n.[C]}{按钮} and a \DV{}{soybean}{}{soy + bean（豆）→ 名词}{n.[C]}{大豆} \NW{}{milk}{}{/mɪlk/}{n.[U]}{奶；（豆）浆} \NW{}{machine}{,}{/məˈʃiːn/}{n.[C]}{机器} an \NW{}{electric}{}{/ɪˈlektrɪk/}{adj.}{电动的；用电的} \NW{}{fan}{}{/fæn/}{n.[C]}{风扇} or an \AG{}{electric}{} \NW{}{bike}{}{/baɪk/}{n.[C]}{自行车} begins to work.\EMK{g}{6} 
\GRAM{6}{\GITEM{动词不定式}{动词不定式在句中可作宾语、表语或状语，常表示目的或将要做的事。}{easy to use 说“用起来”容易；begins to work 中 to work 作 begins 的宾语。}}{}
But a \AG{}{plane}{} on a \AG{}{carrier}{} is \PH{a\ very\ \NW{}{different}{}{/ˈdɪfrənt/}{adj.}{不同的}\ \RV{}{story}{.}}{a different story}{另一回事}

\PARA{6}\SMK{g}An \AG{}{electromagnetic}{} catapult is \RV{}{really}{} a \NW{}{special}{}{/ˈspeʃl/}{adj.}{特殊的} \AG{}{kind}{} of \GL{}{motor}{,}{/ˈməʊtə(r)/}{n.[C]}{电动机} \PH{opened\ out}{open out}{展开} and \NWB{}{laid}{}{lay}{/leɪ/}{vt.}{放置} \NW{}{flat}{.}{/flæt/}{adj.}{平的}\EMK{g}{7} 
\GRAM{7}{\GITEM{分词作定语}{分词可作定语或状语，现在分词表示主动、进行，过去分词表示被动、完成。}{opened out、laid flat 是过去分词，放在 motor 后，说电机被展开、放平。}}{}
\SMK{g}One \NW{}{part}{}{/pɑːt/}{n.[C]}{部分；部件} \NS{}{runs}{}{新义}{run}{/rʌn/}{vi.}{移动；滑行} \RV{}{in}{} a \RV{}{straight}{} \RV{}{line}{} \RV{}{along}{} the other.\EMK{g}{8} 
\GRAM{8}{\GITEM{省略}{前面说过的内容可以省略，或用 so、do、one 代替。}{the other 后省去了 part，指另一个部分，免得重复。}}{}
\SMK{g}\PH{\NW{}{Join}{}{/dʒɔɪn/}{vt.}{连接}\ the\ \NWB{}{moving}{}{move}{/muːv/}{vi.}{移动；活动}\ \AG{}{part}{}\ to}{join A to B}{把A连到B上} a \AG{}{plane}{,} and \PH{there\ you\ have\ it:}{there you have it}{这就成了} a catapult.\EMK{g}{9} 
\GRAM{9}{\GITEM{祈使句}{祈使句用来表示请求、命令或建议，通常省略主语，以动词原形开头。}{Join 原形开头，没有主语，直接告诉读者怎么做。}}{}
\AG{}{Flat}{} motors like this also \NS{}{move}{}{新义}{}{/muːv/}{vt.}{推动；使移动} \GL{}{maglev}{}{/ˈmæɡlev/}{n.[U]}{磁悬浮} \NWB{}{trains}{}{train}{/treɪn/}{n.[C]}{火车} \AG{}{along}{} their \NWB{}{tracks}{.}{track}{/træk/}{n.[C]}{轨道}

\PARA{7}The \NS{}{hard}{}{新词性}{}{/hɑːd/}{adj.}{困难的} \AG{}{part}{} is the \NW{}{power}{}{/ˈpaʊə(r)/}{n.[U]}{电力；动力} \AG{}{system}{.} \NW{}{Firstly}{,}{/ˈfɜːstli/}{adv.}{第一；首先} where does the \AG{}{power}{} come from? A \NW{}{fridge}{}{/frɪdʒ/}{n.[C]}{冰箱} at home \PH{\NWB{}{connects}{}{connect}{/kəˈnekt/}{vi.}{连接}\ to}{connect to}{连接到} the \NW{}{city}{’s}{/ˈsɪti/}{n.[C]}{城市} \AG{}{power}{} \NW{}{network}{,}{/ˈnetwɜːk/}{n.[C]}{网络；电网} but a \AG{}{carrier}{} has to make all of its \RV{}{own}{} \AG{}{electricity}{.} The Fujian makes 190 \GL{}{megawatts}{}{/ˈmeɡəwɒt/}{n.[C]}{兆瓦} \PH{in\ \NW{}{total}{,}{/ˈtəʊtl/}{n.[C]}{总数}}{in total}{总共} but a \NW{}{single}{}{/ˈsɪŋɡl/}{adj.}{单个的} \NS{}{throw}{}{新词性}{}{/θrəʊ/}{n.[C]}{抛；（一次）弹射} of the catapult \NWB{}{needs}{}{need}{/niːd/}{vt.}{需要} 60 to 100 megawatts.

\PARA{8}\NW{}{Secondly}{,}{/ˈsekəndli/}{adv.}{第二；其次} can the \AG{}{ship}{} \AG{}{use}{} all that \AG{}{power}{} \PH{at\ \NW{}{once}{?}{/wʌns/}{adv.}{一次}}{at once}{同时；一下子} \SMK{y}The catapult \NW{}{must}{}{/mʌst/}{modal v.}{必须} \NS{}{reach}{}{新义}{}{/riːtʃ/}{vt.}{达到（某水平）} \NW{}{nearly}{}{/ˈnɪəli/}{adv.}{将近；几乎} 100 megawatts \NW{}{within}{}{/wɪˈðɪn/}{prep.}{在……之内} a \AG{}{few}{} \DV{}{thousandths}{}{thousand + -th → 名词}{n.[C]}{千分之一} of a \RV{}{second}{.}\EMK{y}{1} 
\YI{1}{弹射器必须在几千分之一秒内达到将近 100 兆瓦。}
This \AG{}{power}{} can \NW{}{force}{}{/fɔːs/}{vt.}{用力推动} a \AG{}{plane}{} of over 30 \AG{}{tons}{} \AG{}{along}{} the deck, and the \AG{}{plane}{} \AGB{}{reaches}{}{reach} 280 \NWB{}{kilometres}{}{kilometre}{/ˈkɪləmiːtə(r)/}{n.[C]}{千米；公里} \PH{an\ \RV{}{hour}{}}{an hour}{每小时} \AG{}{within}{} three \AGB{}{seconds}{.}{second} \SMK{g}Only \AG{}{China}{} and \PH{the\ \NW{}{United}{}{/juˈnaɪtɪd/}{adj.}{联合的}\ \NWB{}{States}{}{state}{/steɪt/}{n.[C]}{州}}{the United States}{美国} have \NWB{}{mastered}{}{master}{/ˈmɑːstə(r)/}{vt.}{掌握} the \NW{}{technology}{.}{/tekˈnɒlədʒi/}{n.[C]}{技术}\EMK{g}{10}
\GRAM{10}{\GITEM{现在完成时}{过去发生的事，强调它对现在的影响或经历。}{have mastered 是 have + 过去分词，说两国至今已掌握这项技术。}}{}

\PARA{9}\PN{}{America}{}{/əˈmerɪkə/}{n.}{美国} put it on its \GL{}{nuclear}{-}{/ˈnjuːkliə(r)/}{adj.}{核（动力）的}\NS{}{powered}{}{新词性}{power}{/ˈpaʊə(r)/}{vt.}{给……提供动力} \PN{}{Ford-}{}{/fɔːd/}{n.}{福特（级）}\NW{}{class}{}{/klɑːs/}{n.[C]}{级；类} \AGB{}{carriers}{}{carrier} \PH{back\ in}{back in}{早在（某时）} 2013. But \AG{}{China}{’s} \AG{}{system}{} has \PH{gone\ \RV{}{through}{}}{go through}{经历；经受} over 4,000 \NWB{}{tests}{}{test}{/test/}{n.[C]}{测试；试验} \NW{}{without}{}{/wɪˈðaʊt/}{prep.}{没有} a \AG{}{single}{} \NW{}{accident}{,}{/ˈæksɪdənt/}{n.[C]}{事故} \NW{}{while}{}{/waɪl/}{conj.}{而（表示对比）} the \PN{}{American}{}{/əˈmerɪkən/}{adj.}{美国的} \AG{}{system}{} \PH{\NWB{}{breaks}{}{break}{/breɪk/}{vi.}{坏；出故障}\ down}{break down}{出故障} about \AG{}{once}{} \NS{}{in}{}{新义}{}{/ɪn/}{prep.}{（比率）每……中} \RV{}{every}{} 240 \AGB{}{throws}{.}{throw} \SMK{y}\SMK{g}And the Fujian, with no nuclear \AG{}{power}{,} can do what only a big nuclear \AG{}{carrier}{} \AG{}{could}{} do before.\EMK{y}{2}\EMK{g}{11} 
\YI{2}{而福建舰没有核动力，却能做到以前只有大型核动力航母才能做到的事。}

\GRAM{11}{\GITEM{宾语从句}{宾语从句在复合句中作宾语，从句要用陈述语序。}{what 引出的从句作 do 的宾语，从句里仍是主语在前、谓语 could do 在后。}}{\SEG{the Fujian}{主语（专有名词）}{segA}\SEG{with no nuclear power}{状语（介词短语，说情况）}{segB}\SEG{can do}{谓语（情态动词 + 原形）}{segC}\SEG{what only a big nuclear carrier could do before}{宾语从句（作 do 的宾语）}{segD}}
One \NW{}{reason}{}{/ˈriːzn/}{n.[C]}{原因} is \PH{\NW{}{worth}{}{/wɜːθ/}{adj.}{值得的}\ a\ look.}{be worth a look}{值得一看}

\PARA{10}\SMK{g}Over 100 years \RV{}{ago}{,} \PN{}{AC}{}{/ˌeɪ ˈsiː/}{n.}{交流电} \NW{}{beat}{}{/biːt/}{vt.}{战胜；打败} \PN{}{DC}{}{/ˌdiː ˈsiː/}{n.}{直流电} because it \AG{}{could}{} \RV{}{travel}{} \NW{}{well}{}{/wel/}{adv.}{好；顺利地} over \RV{}{long}{} \NWB{}{distances}{.}{distance}{/ˈdɪstəns/}{n.[C]}{距离}\EMK{g}{12} 
\GRAM{12}{\GITEM{原因状语从句}{原因状语从句说明事情发生的原因，结果状语从句说明由此产生的结果。}{because 从句说交流电获胜的原因：能远距离输送。}}{\SEG{Over 100 years ago}{状语（副词短语，说何时）}{segA}\SEG{AC}{主语（专有名词）}{segB}\SEG{beat}{谓语（动词过去式）}{segC}\SEG{DC}{宾语（专有名词）}{segD}\SEG{because it could travel well over long distances}{原因状语从句（说原因）}{segA}}
But \NS{}{in}{}{新义}{}{/ɪn/}{prep.}{在……方面} many other \NS{}{ways}{}{新义}{way}{/weɪ/}{n.[C]}{方面} DC is \NW{}{ahead}{:}{/əˈhed/}{adv.}{领先；在前} it works \NW{}{fast}{}{/fɑːst/}{adv.}{快地} and gives \NW{}{smooth}{}{/smuːð/}{adj.}{平稳的} \AG{}{power}{.} America started \NS{}{early}{,}{新词性}{}{/ˈɜːli/}{adv.}{早} \NS{}{when}{}{新义}{}{/wen/}{adv.}{（在）那时} DC \AG{}{technology}{} was \NS{}{still}{}{新义}{}{/stɪl/}{adv.}{仍然} \NW{}{weak}{.}{/wiːk/}{adj.}{弱的；差的} \AG{}{China}{} started \NW{}{late}{,}{/leɪt/}{adv.}{晚；迟} but this helped it: it \NWB{}{chose}{}{choose}{/tʃuːz/}{vt.}{选择} a DC \NW{}{electrical}{}{/ɪˈlektrɪkl/}{adj.}{电的；电气的} \AG{}{power}{} \AG{}{system}{} for the \RV{}{whole}{} \AG{}{ship}{.}

\PARA{11}The \RV{}{same}{} \AG{}{system}{} \NW{}{may}{}{/meɪ/}{modal v.}{可能} also \AG{}{power}{} \DV{}{railguns}{,}{rail + gun → 名词}{n.[C]}{电磁炮} \GL{}{laser}{}{/ˈleɪzə(r)/}{n.[C]}{激光} weapons and \GL{}{radars}{.}{/ˈreɪdɑː(r)/}{n.[C]}{雷达} \SMK{g}\NW{}{Chinese}{}{/ˌtʃaɪˈniːz/}{adj.}{中国的} people \NW{}{call}{}{/kɔːl/}{vt.}{把……叫作} them \AG{“}{electric}{} \NWB{}{tigers}{”,}{tiger}{/ˈtaɪɡə(r)/}{n.[C]}{老虎} because they \PH{eat\ up}{eat up}{耗掉；吃光} \RV{}{huge}{} \PH{\NWB{}{amounts}{}{amount}{/əˈmaʊnt/}{n.[C]}{数量}\ of}{amount(s) of}{……的量} \AG{}{electricity}{.}\EMK{g}{13} 
\GRAM{13}{\GITEM{主语 + 谓语 + 宾语 + 宾语补足语}{主语 + 谓语 + 宾语 + 宾语补足语，补足语说明宾语的状态或动作。}{electric tigers 补充说明 them 被叫作什么，去掉它句子就不完整。}}{\SEG{Chinese people}{主语（名词短语）}{segA}\SEG{call}{谓语（动词）}{segB}\SEG{them}{宾语（代词）}{segC}\SEG{“electric tigers”}{宾语补足语（名词短语，说叫作什么）}{segD}\SEG{because they eat up huge amounts of electricity}{原因状语从句（说原因）}{segA}}
\SMK{g}One day, this \AG{}{technology}{} \AG{}{could}{} \AG{}{even}{} help \NWB{}{cars}{}{car}{/kɑː(r)/}{n.[C]}{汽车} \PH{\RV{}{fill}{}\ up\ with}{fill up with}{充满；充上} \AG{}{power}{} \PH{in\ \RVB{}{moments}{.}{moment}}{in moments}{转眼间}\EMK{g}{14}
\GRAM{14}{\GITEM{动词 + 宾语 + do}{want/ask sb to do；let/make/see sb do（不带 to）。}{help cars fill up：help 后先接宾语 cars，再接不带 to 的原形 fill。}}{}
\end{lessonbody}
\sourceline{Adapted from: “福建舰电磁弹射五代机！连美国航母都做不到，到底有多厉害？”, 果壳 (Guokr), 2025-09-27.}
\SelfTest{143}{\TW{September}
\TW{China}
\TW{share}
\TW{news}
\TW{kind}
\TW{plane}
\TW{early}
\TW{warning}
\TW{finish}
\TW{training}
\TW{throw}
\TW{full}
\TW{oil}
\TW{hat}
\TW{off}
\TW{ah}
\TW{exactly}
\TW{such}
\TW{deal}
\TW{airport}
\TW{sea}
\TW{speed}
\TW{rise}
\TW{weigh}
\TW{ton}
\TW{hundred}
\TW{metre}
\TW{enough}
\TW{usual}
\TW{answer}
\TW{pressure}
\TW{forward}
\TW{slow}
\TW{emergency}
\TW{ship}
\TW{bring}
\TW{boil}
\TW{wait}
\TW{could}
\TW{electricity}
\TW{job}
\TW{quickly}
\TW{easy}
\TW{press}
\TW{button}
\TW{milk}
\TW{machine}
\TW{electric}
\TW{fan}
\TW{bike}
\TW{different}
\TW{special}
\TW{lay}
\TW{flat}
\TW{part}
\TW{join}
\TW{move}
\TW{train}
\TW{track}
\TW{power}
\TW{firstly}
\TW{fridge}
\TW{connect}
\TW{city}
\TW{network}
\TW{total}
\TW{single}
\TW{need}
\TW{secondly}
\TW{once}
\TW{must}
\TW{nearly}
\TW{within}
\TW{force}
\TW{kilometre}
\TW{united}
\TW{state}
\TW{master}
\TW{technology}
\TW{class}
\TW{test}
\TW{without}
\TW{accident}
\TW{while}
\TW{break}
\TW{reason}
\TW{worth}
\TW{beat}
\TW{well}
\TW{distance}
\TW{ahead}
\TW{fast}
\TW{smooth}
\TW{weak}
\TW{late}
\TW{choose}
\TW{electrical}
\TW{may}
\TW{Chinese}
\TW{call}
\TW{tiger}
\TW{amount}
\TW{car}
\TW{air\,\nstag{新义}}
\TW{sound\,\nstag{新词性}}
\TW{use\,\nstag{新词性}}
\TW{run\,\nstag{新义}}
\TW{move\,\nstag{新义}}
\TW{hard\,\nstag{新词性}}
\TW{throw\,\nstag{新词性}}
\TW{reach\,\nstag{新义}}
\TW{power\,\nstag{新词性}}
\TW{in\,\nstag{新义}}
\TW{way\,\nstag{新义}}
\TW{early\,\nstag{新词性}}
\TW{when\,\nstag{新义}}
\TW{still\,\nstag{新义}}
\TW{kind(s) of}
\TW{be full of}
\TW{hats off (to sb)}
\TW{a big deal}
\TW{at sea}
\TW{speed up}
\TW{tens of}
\TW{in an emergency}
\TW{bring sth to the boil}
\TW{a different story}
\TW{open out}
\TW{join A to B}
\TW{there you have it}
\TW{connect to}
\TW{in total}
\TW{at once}
\TW{an hour}
\TW{the United States}
\TW{back in}
\TW{go through}
\TW{break down}
\TW{be worth a look}
\TW{eat up}
\TW{amount(s) of}
\TW{fill up with}
\TW{in moments}}
\par\write-1{LESSONEND 2 \thepage}
